\chapter{Моделирование многозадачной координации в группе роботизированных НТТС}\label{ch:ch2}

Распределение задач между членами группы исполнителей -- классическая проблема с богатой историей и множеством применений. Традиционно распределение задач состоит в назначении каждому исполнителю одной из многих возможных задач таким образом, чтобы была оптимизирована метрика производительности и эффективности. Однако, в зависимости от структуры целевой функции, эта проблема может быть трудоемкой с точки зрения вычислений, т.\ к. требуется \textit{единый} центр принятия решений. Такая реализация требует устойчивой связи между НТТС, что может быть сложно реализовать в крупномасштабных системах из-за помех. Альтернативой этому подходу является предоставление НТТС возможности определять задачи и позволять им самостоятельно принимать решения о том, какую работу выполнять на основе общей стратегии в рамках своих ресурсов.

При этом возникает такой конфликт, что несколько НТТС в группе могут претендовать на выполнение одной и той же задачи одновременно. Это не будет являться эффективной стратегией. Для разрешения таких столкновений интересов и определения рациональной общей стратегии принятия решения предлагается воспользоваться математическими методами Теории игр \cite{GlobalGames, GlobalGames2, GlobalGamesRus, GlobalGamesRus2}. При таком подходе состояние системы будет частично контролироваться НТТС, которые имеют минимальное взаимодействие между собой и чьи действия будут влиять на других членов группы в пределах их зоны контроля. Преимущество заключается в том, что при простых предположениях о структуре функции вознаграждения НТТС в системе будут присутствовать байесовские равновесия Нэша (БРН).

\nomenclature{БРН}{байесовские равновесия Нэша\nomrefpage}

Графический замысел системы изображен на рис.\ \cref{fig:system_model}. По нему, конечное число НТТС должно выбирать из группового массива задачи. Каждый НТТС может иметь разное навесное оборудование для выполнения своего типа задач, а также своей ограниченный ресурс. Очередь пополняется задачами из внешнего источника и от НТТС по распределению Пуассона. Каждая такая задача характеризуется независимыми случайными величинами, соответствующими их уровню сложности (сколько ресурсов будет затрачено на ее выполнение) и типом требуемого оборудования для ее выполнения. Задачи распределяются равномерно между всеми НТТС в группе. Координация между НТТС, необходимая для выполнения общей задачи, происходит посредством локальных взаимодействий по сети. Наконец, НТТС решает, какую из задач выполнить, и связанный с этим выигрыш зависит от выбранной задачи и от количества НТТС, принявших такое же решение.

\begin{figure}[ht]
	\centering{
		\includegraphics[scale=0.4]{drawio/system-math.drawio.png}
	}
	\caption{Графический замысел модели системы}\label{fig:system_model}
\end{figure}

\section{Описание модели системы}\label{sec:ch2/sec2}

Рассмотрим систему с $ N>1 $ НТТС входящими в множество $ i \in \{ 1; \dotsc; N \} = [N] $. НТТС являются автономными роботизированными наземными транспортно-технологическими средствами, принимающими решения и выбирающими одну из задач групповой очереди, над которыми они должны работать. Обозначим через $R_i$ $i$-й НТТС. Каждый $R_i$ характеризуется набором характеристик $\langle \text{Eqp}_i, E_i \rangle$, включающим тип оборудования и объем доступного ресурса соответственно.

Поступление новых задач в систему осуществляется согласно двум независимым пуассоновским процессам: внешнему потоку $\Pi_s(t)$ с интенсивностью $\lambda_s$ и локальному потоку $\Pi_l(t)$ с интенсивностью $\lambda_l$.

Интенсивности $\lambda_s$ и $\lambda_l$ характеризуют среднее число задач, поступающих в единицу времени от соответствующих источников.

Процесс $\Pi(t)$, порождающий входящие задачи, является пуассоновским, если для любого временного интервала длины $t$ вероятность поступления ровно $k$ задач определяется законом Пуассона:

\begin{equation}
	P_k(t) = \Pr\left(\Pi(t) = k\right) = \frac{(\lambda t)^k}{k!} e^{-\lambda t}, \quad k = 0, 1, 2, \dotsc
\end{equation}

где $\lambda$ — интенсивность процесса.

Для двух независимых пуассоновских процессов суммарный входящий поток также является пуассоновским процессом с интенсивностью $\lambda = \lambda_s + \lambda_l$. Таким образом, групповую очередь задач, формируемую как внешним источником, так и НТТС, следует рассматривать, как единый поток задач для всей группы НТТС.

Каждой $j$-й задаче $T_j$ соответствует набор характеристик $\langle \Theta_{T_j}, \text{Req}_j \rangle$ включающий сложность выполнения (количество затрачиваемых ресурсов) задачи и~необходимый тип оборудования для ее выполнения.

$ \Theta_T $ выбирается из гауссова распределения:

\begin{equation}
	\Theta_T \sim \mathcal{N}(0; \sigma_T^2)
\end{equation}

Для корректности модели, учитывая, что $\Theta_T$ не может быть отрицательным, будем использовать усечённое нормальное распределение:

\begin{equation}
	\Theta_T \sim \mathcal{N}^+(0; \sigma_T^2),
\end{equation}

где верхний индекс $+$ означает усечение на уровне нуля.

У каждого $R_i$ есть ограниченный запас энергии $E_i \in [0, E_{\text{max}}]$, где $E_{\text{max}}$ -- максимальный запас энергии. Выполнение задачи $ T $ требует затрат $ e_T > 0 $.

Динамика изменения энергии для $R_i$ описывается следующим образом:

\begin{equation}
	E_i^{(T+1)} = E_i^{(T)} - e_T \cdot \mathbf{1}(A_i = T) + \delta \cdot \Delta t,
\end{equation}

где $\delta > 0$ -- коэффициент восстановления энергии, $\Delta t$ -- время между обновлениями, a $\mathbf{1}(A_i = T)$ — индикаторная функция, принимающая значение $1$, если $R_i$ выбрал задачу $T$, и $0$ в противном случае.

Общий групповой массив задач в системе задается множеством:

\begin{equation}
	\mathcal{T}(t) = \{ T_1, T_2, \dotsc, T_M(t) \},
\end{equation}

где $M(t)$ это текущее количество доступных задач в момент времени $t$.

Каждый $R_i$ способен автономно принимать решения и выбирать задачи из $\mathcal{T}(t)$.

Выбор задачи каждым $R_i$ ограничен. Задача $T_j \in \mathcal{T}(t)$ может быть назначена $R_i$, только если одновременно выполняются следующие два условия:
\begin{enumerate}
	\item Совместимость оборудования: $\text{Eqp}_i \succeq \text{Req}_{T_j}$,
	\item Достаточность ресурса: $ E_i \geq \Theta_{T_j} $,
\end{enumerate}
где отношение $\succeq$ означает, что оборудование $R_i$ удовлетворяет или превосходит требования задачи $T_j$, а $\Theta_{T_j}$~— объем ресурса, необходимого для выполнения задачи $T_j$, т.\ е.:

\begin{equation}
	T_j \in \mathcal{T} \text{ доступна для } R_i \iff
	\begin{cases}
		\text{Eqp}_i \succeq \text{Req}_{T_j}, \\
		E_i \geq \Theta_{T_j}.
	\end{cases}
\end{equation}

Матрица соответствия задач для $R_i$ будет определена как $A^{(i)} \in \{0,1\}^{M(t) \times 1}$, где каждый элемент $a_j^{(i)}$ указывает, доступна ли задача $T_j$ для $R_i$:

\begin{equation}
	a_j^{(i)} =
	\begin{cases}
		1, & \text{если } \text{Eqp}_i \succeq \text{Req}_{T_j} \text{ и } E_i \geq \Theta_{T_j}, \\
		0, & \text{иначе}.
	\end{cases}
\end{equation}

Таким образом, $B^{(i)} = (b_1^{(i)}, b_2^{(i)}, \dotsc, b_{M(t)}^{(i)})^\top$ описывает, что для каждого $R_i$ есть локальный массив задач из общего группового массива задач. При этом, одна и та же задача может быть доступна разным $R$.

Предположим, что $R_i$ взаимодействуют в соответствии с графом влияния $ \mathcal{G} = ([N],\mathcal{E}) $, где $\mathcal{E}$ множество ребер графа. А также, что граф связан, неориентирован и имеет единичные веса. Пусть $ \mathcal{R}_i $ обозначает множество соседей $R_i$, $ i \in [N] $.

Рассмотрим простейший случай, когда имеются два $R i, j \in [N] $ и $ (i,j) \in \mathcal{E} $, а также две задачи $T$ такие, что $[T_1, T_2] \in B^{(1)}$ и $[T_1, T_2] \in B^{(2)}$.

Так, каждый $R$ принимает двоичное решение $ A_i \in \{ 1,2 \} $, которое соответствует $R$, решающему выполнить $T_1$ или $ T_2 $, соответственно. $R_i$ использует измеримое правило $ \gamma_i: \mathbb{R}^2 \rightarrow \{ 1,2 \} $ по своим наблюдениям, чтобы вычислить требуемое решение. Правило представляет собой отображение из пространства наблюдения сигнала о решении для выполнения $T_1$ или $ T_2$. Пусть $ \Upsilon^{(i)} = [ \Upsilon_1^{(i)}, \Upsilon_2^{(i)} ]^{\mathsf{T}} $, тогда действие $R_i$ вычисляется в соответствии с:

\begin{equation}
	\label{eq:equation4}
	A_i = \gamma_i\left( \Upsilon^{\left(i\right)} \right).
\end{equation}

Поскольку $R$ связаны, можно предположить, что они участвуют в \textbf{координационной игре}, определяемой матрицей взаимодействия (табл.~\cref{tab:binGameTable}). Идея, лежащая в основе матрицы взаимодействия, заключается в том, что задача с затратами $ \theta_T $ становится легче по мере увеличения количества $R$ в системе. Однако задача будет успешно выполнена только в том случае, если достаточная доля $R$ в~окружении $ \mathcal{R}_i $ будет сотрудничать в выполнении одной и той же задачи. То есть выигрыш каждого $R$ зависит от выбора $ T $ и $ \mathcal{R}_i $.

\begin{table}[htbp]
	\captionsetup[table]{format=split}
	\centering
	% Увеличиваем высоту строк в таблице
	\renewcommand{\arraystretch}{1.5} % Можно изменять значение для более крупного размера
	\caption{Матрица бинарной игры между двумя НТТС}%
	\label{tab:binGameTable}%
	\begin{tabular}{c c|c|c|}
		\multicolumn{2}{c}{} & \multicolumn{2}{c}{$R_j$} \\ % Заголовок без рамок
		\multicolumn{2}{c}{} & \multicolumn{1}{c}{1} & \multicolumn{1}{c}{2} \\ \cline{3-4}
		\multirow{2}{*}{$R_i$} & 1 & $\left( \frac{1}{|\mathcal{N}_i|} - \frac{\theta_1}{N}, \frac{1}{|\mathcal{N}_j|} - \frac{\theta_1}{N} \right)$ & $\left( -\frac{\theta_1}{N}, -\frac{\theta_2}{N} \right)$ \\ \cline{3-4}
		& 2 & $\left( -\frac{\theta_2}{N}, -\frac{\theta_1}{N} \right)$ & $\left( \frac{1}{|\mathcal{N}_i|} - \frac{\theta_2}{N}, \frac{1}{|\mathcal{N}_j|} - \frac{\theta_2}{N} \right)$ \\ \cline{3-4}
	\end{tabular}
\end{table}

\begin{subequations}
	Рассмотрим матрицу бинарной игры (табл.\ \cref{tab:binGameTable}). Если два НТТС выбирают одну и ту же задачу, их выигрыш увеличивается, так как они координируют свои усилия. Если они выбирают разные задачи, их выигрыш уменьшается.

	Пусть $ L_{ij} \equiv \sfrac{N}{min\{|\mathcal{N}_i|,|\mathcal{N}_j|\}} $ и $ \mathcal{S}_{ij} $ обозначает множество БРН с чистой стратегией детерминированной координационной игры:
	\label{eq:equation5}
	\begin{align}
		\theta_1 - \theta_2 < -L_{ij} 					& \implies \mathcal{S}_{ij}=\{(1,1)\}; \label{eq:equation5_1} \\
		-L_{ij} \leq \theta_1 - \theta_2 \leq L_{ij} 	& \implies \mathcal{S}_{ij}=\{(1,1),(2,2)\}; \label{eq:equation5_2} \\
		\theta_1 - \theta_2 > L_{ij} 					& \implies \mathcal{S}_{ij}=\{(2,2)\}. \label{eq:equation5_3}
	\end{align}
\end{subequations}

Утверждение проверяется с использованием определения БРН \cite{GlobalGames3}. Т.\ о. в описанной игре существует несколько БРН, которые зависят от разности сложностей задач $ \theta $ и структуры графа взаимодействия $ \mathcal{G} $.

Следовательно, существует такой режим работы, в котором НТТС $ i $ и $ j $ могут выбрать одно из нескольких равновесий в зависимости от значений переменных $ \theta_1 $ и $ \theta_2 $. Например, если верно выражение \eqref{eq:equation5_1}, то все НТТС выбирают задачу $ \theta_1 $. Если верно \eqref{eq:equation5_2}, то возможны два равновесия, когда все выбирают задачу $ \theta_1 $ или $ \theta_2 $. А при \eqref{eq:equation5_3} -- все выбирают задачу $ \theta_2 $ соответственно. Более того, набор равновесий зависит от структуры окрестностей графа.

Пусть $ \theta \equiv [\theta_1\theta_2]^{\mathtt{T}} $. Полезность НТТС $ i $ в игре с одним из его соседей $ j $ тогда определяется:

\begin{equation}
	\label{eq:equation6}
	\begin{split}
	\mathcal{U}(a_i, a_j, \theta) \equiv \textbf{1}(a_i=1)\left(\frac{1}{|\mathcal{N}_i|}\textbf{1}(a_j=1)-\frac{\theta_1}{N}\right)+ \\
	+\textbf{1}(a_i=2)\left(\frac{1}{|\mathcal{N}_j|}\textbf{1}(a_j=2)-\frac{\theta_2}{N}\right).
	\end{split}
\end{equation}

В координационной игре НТТС $ i $ одновременно выполняет одно действие со всеми своими соседями $ j $ в $ N_i $ \cite{GlobalGames4}. Следовательно, можно определить следующую локальную функцию полезности:

\begin{equation}
	\label{eq:equation7}
	\mathcal{V}_i(a_i, a_j, \theta) \equiv
	\displaystyle\sum_{j \in N_i} \mathcal{U}_i(a_i, a_j, \theta),
\end{equation}
где полезность НТТС $ i $ определяется как сумма выигрышей от взаимодействия со всеми соседями.

НТТС в такой игре действуют не кооперативным образом, стремясь получить свой индивидуальный максимальный выигрыш в отношении своей стратегической политики \cite{GlobalGames3}. Другими словами, НТТС стремятся увеличить свою локальную полезность, выбирая стратегию $ \gamma_i $, которая повышает ожидаемый выигрыш.

Пусть $ \Gamma $ -- профиль стратегии координации, т.\ е. набор политик координации, используемых всеми НТТС в системе. Тогда
\begin{equation}
	\label{eq:equation8}
	\Gamma \equiv (\gamma_1, \dots, \gamma_N),
\end{equation}
где $ \gamma_i: \mathbb{R}^2 \rightarrow \{1,2\} $ является измеримой функцией, $ i \in [N] $. Аналогично, пусть $ \Gamma_{-i} $ обозначает профиль стратегии координации, используемый всеми НТТС, кроме $ i $.

Учитывая $ \Gamma_{-i} $, целью $ i $-го НТТС является решение следующей (стохастической) задачи оптимизации:

\begin{equation}
	\label{eq:equation9}
	\max_{\gamma_i}\mathcal{J}_i(\gamma_i, \Gamma_{-i}) \equiv \textbf{E}\left[\mathcal{V}_i(A_i, A_{-i}, \Theta)\right],
\end{equation}
где математическое ожидание берется для экзогенных случайных величин $\Theta$ и $ Z_t^i, t \in \{1,2\} $.

Таким образом, стратегия для профиля $ \gamma^*=(\gamma_i^*, \dots, \gamma_N^*) $ является БРН для многозадачной гауссовой глобальной игры, если:

\begin{equation}
	\label{eq:equation10}
	\mathcal{J}_i\left(\gamma_i^*,\gamma_{-i}^*\right) \geq \mathcal{J}_i\left(\gamma_i,\gamma_{-i}^*\right), \gamma_i \in \Gamma_i, \forall i \in [N],
\end{equation}
т.\ е. стратегия $ \gamma^*=(\gamma_i^*, \dots, \gamma_N^*) $ является БРН, если ни одно НТТС не может увеличить свой ожидаемый выигрыш, отклоняясь от своей стратегии при фиксированных стратегиях других НТТС.


\fixme{Построить обобщённый пример  игры, где НТТС выбирают задачи из массива, чтобы максимизировать свою долю полезности, учитывая неопределённость в действиях других. Найти Байесовское равновесие Нэша (БРН) — набор стратегий, при котором ни один агент не хочет отклониться, если он знает лишь вероятностное распределение поведения других. Построить матрицу игры. Найти графически и аналитически оптимальную (смешанную) стратегию, цену игры V, нижнюю α и верхнюю β цены игры в чистых стратегиях}


\section{Стратегии выбора задач}

НТТС могут использовать следующие стратегии для выбора задач:

\begin{enumerate}
    \item \textbf{Энергосберегающая стратегия}: НТТС выбирает задачу с минимальными затратами энергии $e_t$ при условии $E_i \geq e_t$.

    \item \textbf{Стратегия минимизации задержек}: НТТС выбирает задачу, которая дольше всего находится в очереди.

    \item \textbf{Стратегия баланса}: НТТС выбирает задачу, максимизирующую отношение выигрыша к затратам энергии.
\end{enumerate}

\section*{Выводы по главе}\label{sec:ch2/sec5}

Выводы

\FloatBarrier
